% !TeX root = ../thesis.tex

\chapter*{Zusammenfassung}
\addcontentsline{toc}{chapter}{Zusammenfassung}

Moderne Sequenzierungstechnologien ermöglichen es, mikrobielle Gemeinschaften in einem bislang nicht erreichten Umfang zu charakterisieren. Die dabei entstehenden Daten stellen die statistische Analyse jedoch vor erhebliche Herausforderungen. Mikrobiomdaten sind hochdimensional, spärlich besetzt und kompositionell, und die resultierenden Schlussfolgerungen können stark von Entscheidungen in der Datenaufbereitung und Modellierung abhängen. Diese Dissertation adressiert diese Herausforderungen durch zwei methodische Beiträge: einen umfassenden Workflow zur Analyse mikrobieller Assoziationsnetzwerke, sowie ein Verfahren zur Verbesserung der Approximation von $p$-Werten bei Permutationstests, wenn nur eine begrenzte Anzahl von Permutationen durchführbar ist. Beide Beiträge werden in eine übergreifende Perspektive eingebettet, in der Analysen auf Community-Ebene, Analysen auf Feature-Ebene, und Netzwerkanalysen komplementäre Sichtweisen auf dieselben Mikrobiomdaten liefern.

Der erste Hauptbeitrag ist \texttt{NetCoMi}, ein R-Paket zur Konstruktion, Analyse, Visualisierung und zum Vergleich mikrobieller Assoziationsnetzwerke. Das Lernen solcher Netzwerke erfordert eine Reihe miteinander verknüpfter Entscheidungen hinsichtlich des Filterns von Stichproben und Taxa, Umgang mit Nullen, Datentransformation, Assoziationsschätzung und Sparsifizierung. \texttt{NetCoMi} integriert diese Schritte in einen kohärenten Workflow und erlaubt es zugleich, zwischen etablierten Methoden zu wählen oder zuvor erstellte Assoziationsmatrizen zu übergeben. Das Framework unterstützt darüber hinaus die Charakterisierung der erstellen Netzwerke durch globale Eigenschaften, Zentralitätsmaße, Clusterverfahren und graphlet-basierte Maße, sowie vergleichende Analysen gruppenspezifischer Netzwerke. Die Dissertation behandelt sowohl die methodischen Grundlagen dieser Komponenten als auch die Implementierungsherausforderungen, die sich aus der Kombination von Methoden mit unterschiedlichen Annahmen und Datenanforderungen ergeben.

Der zweite Hauptbeitrag ist \texttt{permApprox}, ein Framework zur Verfeinerung permutationsbasierter $p$-Werte durch parametrische Approximation des relevanten Randbereichs der Permutationsverteilung. Permutationstests sind in der Mikrobiomanalyse besonders nützlich, wenn für komplexe Teststatistiken keine verlässliche parametrische Nullverteilung verfügbar ist. Die begrenzte Zahl durchführbarer Permutationen führt jedoch häufig zu einer groben Auflösung empirischer $p$-Werte, was insbesondere bei vielen simultanen Tests problematisch ist. Ansätze auf Basis der verallgemeinerten Pareto-Verteilung können diese Auflösung verbessern, aber bei beschränktem Träger zu $p$-Werten von exakt null führen. \texttt{permApprox} vermeidet dies durch Parameterschätzung unter Nebenbedingungen und gewährleistet dadurch strikt positive $p$-Werte. Neben der eigentlichen Methode, bietet \texttt{permApprox} einen kompletten Workflow, welcher weitere Schritte des Approximationsverfahrens beinhaltet, wie das Finden eines Schwellwertes für den Randbereich und eine valide Parameterschätzung. Simulationsstudien zeigen die Genauigkeit, Robustheit und das Vermeiden von Null-$p$-Werten bei begrenzten Permutationsbudgets.

Zur Veranschaulichung der methodischen Beiträge wird eine integrierte Analyse von Sequenzierungsdaten zum Darmmikrobiom von Säuglingen aus der PASTURE-Kohorte durchgeführt. Sie umfasst Analysen auf Community-Ebene, Feature-basierte Vergleiche einzelner Taxa zwischen Gruppen unterschiedlicher Dauer exklusiven Stillens, sowie die Konstruktion und den Vergleich mikrobieller Assoziationsnetzwerke mit \texttt{NetCoMi}. Für alle Schritte werden permutationsbasierte Verfahren eingesetzt, sodass die resultierenden $p$-Werte mit \texttt{permApprox} verfeinert werden können. Die Anwendung zeigt, dass unterschiedliche Analyseziele nicht zwangsläufig dieselben Gattungen hervorheben: Taxa mit marginalen Häufigkeitsunterschieden sind nicht unbedingt zentrale Knoten in den Netzwerken, und differentielle Assoziationen können ohne ausgeprägte marginale Gruppenunterschiede auftreten. Netzwerkanalysen ergänzen andere Ansätze der Mikrobiomanalyse daher, anstatt sie zu ersetzen.

Insgesamt argumentiert diese Dissertation, dass verlässliche mikrobiomstatistische Inferenz vom vollständigen Analyseworkflow und nicht von einer einzelnen Methode oder Teststatistik abhängt. Entscheidungen in der Datenaufbereitung definieren wesentliche Charakteristika der Daten und können alle darauf aufbauenden Ergebnisse beeinflussen. Methodische Flexibilität ist wertvoll, um Analysen an konkrete Datensätze anzupassen, erfordert jedoch transparente Berichterstattung, Sensitivitätsanalysen und eine vorsichtige Interpretation. Die in dieser Dissertation entwickelten Methoden und Softwarewerkzeuge sollen komplexe Mikrobiomanalysen zugänglicher, reproduzierbarer und statistisch besser begründbar machen, ohne die Unsicherheit zu verschleiern, die in spärlich besetzten und kompositionellen Sequenzierungsdaten inhärent ist.


\chapter*{Summary}
\addcontentsline{toc}{chapter}{Summary}

High-throughput sequencing has made it possible to characterize microbial communities at an unprecedented scale, but the resulting data pose substantial statistical challenges. Microbiome count data are high-dimensional, sparse, and compositional, and conclusions may depend strongly on choices made during preprocessing and modeling. This thesis addresses these challenges through two methodological contributions: an integrated framework for microbial association network analysis and a framework for improving permutation-based p-value approximation when only a limited number of permutations is feasible. These contributions are embedded in a broader perspective in which community-level summaries, feature-level analyses, and association networks provide complementary views on the same microbiome data.

The first main contribution is \texttt{NetCoMi}, an R package for constructing, analyzing, visualizing, and comparing microbial association networks. Network learning requires a sequence of interconnected decisions concerning sample and taxa filtering, zero handling, data transformation, association estimation, and sparsification. \texttt{NetCoMi} integrates these steps into a coherent workflow while allowing users to select among established methods or supply externally estimated association matrices. The framework further supports the characterization of network structure through global properties, centrality measures, clustering, and graphlet-based summaries, as well as comparative analyses of group-specific networks. The thesis discusses both the methodological foundations of these components and the implementation challenges involved in combining methods with different assumptions and data requirements.

The second main contribution is \texttt{permApprox}, a general framework for refining permutation p-values by approximating the relevant tail of the permutation distribution with a parametric model. Permutation tests are particularly useful in microbiome analysis because many statistics arising after complex preprocessing or modelling lack tractable parametric null distributions. However, computational constraints often limit the number of feasible permutations, resulting in coarse empirical p-value resolution. This is especially problematic in settings involving a large number of simultaneous tests. Existing approximation approaches based on the generalized Pareto distribution can improve this resolution, but may yield zero p-values when the fitted tail has bounded support. \texttt{permApprox} addresses this issue through constrained fitting, ensuring strictly positive p-value estimates, while also providing procedures for further steps of the approximation workflow, including threshold selection and valid parameter estimation. Simulation studies demonstrate the accuracy and robustness of the proposed approach, as well as its ability to avoid zero p-values under limited permutation budgets.

The methodological developments are illustrated in an integrated analysis of infant gut microbiome data from the PASTURE cohort. It comprises community-level analyses, feature-level comparisons of individual taxa between groups defined by the duration of exclusive breastfeeding, and the construction and comparison of microbial association networks using \texttt{NetCoMi}. Permutation-based procedures are used throughout so that the resulting $p$-values can be refined with \texttt{permApprox}. The application shows that different analytical targets do not necessarily highlight the same genera: taxa with marginal abundance differences are not necessarily central nodes in the networks, and differential associations may occur without pronounced marginal group differences. Network analyses therefore complement, rather than replace, other approaches to microbiome analysis.

Overall, this thesis argues that reliable microbiome inference depends on the complete analytical workflow rather than on an individual method or test statistic. Preprocessing choices define essential data characteristics and can affect all subsequent results. Methodological flexibility is valuable for adapting analyses to specific data sets, but it also requires transparent reporting, sensitivity analyses, and cautious interpretation. The software and methods developed in this thesis aim to make complex microbiome analyses more accessible, reproducible, and statistically defensible while acknowledging the uncertainty inherent in sparse, compositional sequencing data.


\cleardoublepage